\documentclass[11pt]{article}

% Silence some warnings
\hfuzz=5pt
\usepackage{silence}
\WarningFilter{latexfont}{Font shape}
\WarningFilter{latexfont}{Some font shapes were}

% Packages
\usepackage{alt}
\usepackage{kotex}
\usepackage{graphicx}
\usepackage{indentfirst}
\usepackage{microtype}
\usepackage{reactcode}
\usepackage{setspace}

\usepackage{amssymb,amsmath,mathtools} % Before unicode-math
\usepackage[
  math-style=ISO,
  warnings-off={mathtools-colon, mathtools-overbracket},
]{unicode-math}

% Font steup
\usepackage[tt=false]{libertine}
\setmathfont{LibertinusMath-Regular.otf}[Scale=MatchUppercase]
\setmonofont{inconsolata}[Scale=MatchLowercase, StylisticSet=3]
\setmainhangulfont{Noto Serif KR}[
  Scale=MatchUppercase,
  AutoFakeSlant=0.15
]
\setsanshangulfont{KoPubWorldDotum_Pro}[
  Scale=MatchUppercase,
  BoldFont={* Bold}
]

% Temporary; for placeholder
\usepackage{jiwonlipsum}

% Load ACL style
\usepackage[preprint]{acl}
\renewcommand{\abstractname}{초록}
\renewcommand{\citet}[1]{\citeauthor{#1}(\citeyear{#1})}
\ExplSyntaxOn
\NewDocumentCommand \Overline { o m }
{
  \int_compare:nNnTF
  { \str_count:n { #2 } } < 2
  { \overline }
  { \overline }
  {
    \IfNoValueTF { #1 }
    { #2 }
    { \vphantom{#1}\smash{#2} }
  }
}
\ExplSyntaxOff
\onehalfspacing

% Metadata
\title{배열 완성 과제에서 코드 생성 모델의 선호와 예측 단서 분석: \\
리액트 훅의 함수--배열 구조를 중심으로}
\author{박준영 \\
  인문대학 언어학과 \\
  서울대학교 \\
\texttt{bloomwayz@snu.ac.kr} \\}

\begin{document}

\maketitle
\begin{abstract}
  \jiwon[1]
\end{abstract}

\section{서론}
\label{sec:intro}

\alt{대형 언어 모델}{large language models (LLMs)}의 프로그래밍 성능은 계속해서 향상되어 왔으며, 최근 모델은 코드 생성 작업에서 강한 성능을 보인다\citep{austin2021synthesis,fried2023incoder}.
이러한 성능에 힘입어 인공지능 도구는 이미 개발 현장에 폭넓게 자리 잡아 대부분의 개발자가 인공지능 도구를 개발에 활용하고 있고, 그중 많은 이들이 인공지능 도구의 성능에 호의적이라고 밝혔다\citep{Sta25Stack}.

이 추세를 따라 코드 생성 모델의 행동과 내부 구조에 대한 연구도 다방면으로 이루어졌다.
대표적으로 모델이 코드의 문법 구조, 식별자, 데이터 흐름, 제어 흐름, 기능적 동등성과 같은 정보를 어느 정도 표상하는지 분석한 연구가 있고\citep{troshin2022probing,ma2024unveiling,maveli2025functional}, 코드 모델의 예측이 변수명이나 의미 보존 변형과 같은 표면적 변화에 민감할 수 있음을 보인 연구도 있다\citep{yefet2020adversarial,wang2024naming,orvalho2025robust}.
이러한 연구들은 코드 모델의 예측이 프로그램 의미뿐 아니라 다양한 표면적·구조적 단서의 영향을 받을 수 있음을 시사한다.

다만 이들 연구는 주로 완성된 코드의 정답 여부나 오류 유형, 의미 보존 변형에 대한 강건성, 모델 표상의 전반적인 성질을 분석하는 데 초점을 두었다.
특정 위치에서 후보 토큰의 로짓이 어떤 단서에 따라 달라지는지는 아직 충분히 분석되지 않았으며, 특히 라이브러리의 의미 규약과 일반적인 문법 패턴이 함께 작동하는 상황에서 모델이 다음 토큰을 예측할 때 어떤 단서에 기대는지를 국소적으로 살펴본 연구는 찾아보기 어렵다.

이에, 이 글은 대형 언어 모델이 코드의 다음 토큰을 예측할 때 어떤 토큰을 선호하는지 관찰하고, 이때 어떤 정보가 어느 정도로 코드 생성에 영향을 미치는지 분석한다.
특히 이 연구에서는 \alt{리액트}{React} 라이브러리에서 제공되는 함수 가운데 하나인 \texttt{useEffect}의 구조와 유사한 함수--배열 패턴에 모델이 어떻게 반응하는지를 집중적으로 살펴본다.
아래는 우리가 연구에서 사용한 함수--배열 패턴의 예시이다.

% tex-fmt: off
\begin{reactcode}[label=]
const [s, setS] = useState(0);
const c = 0;

useEffect(() => {
  fetch(`/api/me?x=${s}&y=${c}`)
}, [s])
\end{reactcode}
% tex-fmt: on

우선 위 코드에 등장하는 변수를 살펴보자.
첫 번째 줄에서는 상태 변수 \texttt{s}와 상태를 갱신하는 함수 \texttt{setS}를 선언하고 있고, 두 번째 줄에서는 상수 \texttt{c}를 선언하고 있다.
상태 \texttt{s}와 상수 \texttt{c}는 얼핏 닮아 보이지만, 상태는 갱신 함수에 의해 그 값이 달라질 수 있고, 상수의 값은 프로그램이 종료될 때까지 값이 달라지지 않는다는 점에서 둘은 대별된다.

그 다음 집중할 것은 \texttt{useEffect} 함수를 호출하는 대목이다.
\texttt{useEffect}는 또 다른 함수 $f$와 배열 $[\Overline{e}]$를 인자로 받는 함수인데, 이때 배열은 상태 등의 \alt{반응형}{reactive}~값만으로 구성될 것이 강력히 권고된다\citep{met2026sep}.
상태 \texttt{s}와 상수 \texttt{c}는 모두 함수의 본문에 등장하지만 배열에는 \texttt{s}만이 존재하는데, 이는 \texttt{s}는 상태 변수로서 반응형 값이고 \texttt{c}는 상수로서 비반응형 값이기 때문이다.

여기서 몇 가지 의문이 남는다.
대형 언어 모델은 함수--배열 구조에서 배열에 어떤 항목이 올지 올바르게 예측할 수 있을까?
올바르게 예측할 수 있다면, 혹은 예측이 틀리다면, 모델은 어떤 정보를 근거로 삼아 토큰을 예측했을까?
모델이 리액트 라이브러리를 사용하여 코드가 작성되었다는 점을 인식하고 리액트 훅의 의미를 계산했기 때문일까, 아니면 함수와 배열이 매개변수로 쓰이는 문법 구조 때문일까?
그것도 아니라면, \texttt{useState} 같은 키워드를 단서로 삼는 것일까?

이 질문에 답하기 위해, 우리는 45개의 변수 쌍을 바탕으로 자바스크립트 코드 데이터셋을 구축하였다.
각 변수 쌍에 대해 동일한 코드 골격을 유지한 채 함수 이름, 변수 선언 형태, 변수 선언 순서, 변수 이름을 조작하여 40개의 변형 코드를 만들었다.
이 데이터를 이용하여 Llama 3.2 1B 모델과 Gemma 3 1B 모델이 배열의 첫 항목을 예측하도록 했고, 이때 예측되는 토큰의 로짓을 측정하였다.
이어서 함수 본문에서의 변수 사용과 배열을 도입하는 문맥을 각각 조작하는 두 건의 \alt{소거}{ablation}~실험을 통해 관찰된 선호가 무엇으로 환원되는지를 분해하였다.

실험 결과, 우리는 다음 세 가지를 발견할 수 있었다.
\begin{itemize}
  \item \textbf{상태 변수 선호.}
    모델은 실험 조건 전반에서 배열의 첫 요소로서 상태 변수를 선호하는 모습을 보였다.
    이 경향은 리액트 훅 \texttt{useEffect} 또는 \texttt{useLayoutEffect}가 사용되었는지, 일반 자바스크립트 함수 \texttt{subscribe}가 사용되었는지에 상관없이 일관되게 나타나며, 오히려 리액트 환경일 때보다 일반 자바스크립트 환경일 때 모델이 상태 변수를 선호하는 정도가 더 높았다.
  \item \textbf{선언 형태와 변수 사용 단서.}
    이러한 모델의 선호는 두 가지 단서가 더해진 결과로 보인다.
    하나는 변수가 \texttt{useState} 형태로 선언되었다는 `선언 형태' 단서이고, 다른 하나는 그 변수가 함수 본문에서 실제로 쓰였다는 `본문 사용' 단서이다.
    두 단서는 대체로 합쳐지는 방향으로 작동하지만, 본문에서 상태 변수가 쓰이지 않은 상황에서 어느 단서가 더 강하게 작용하는지는 배열을 도입하는 문맥이 리액트 환경인지의 여부에 따라 달라진다.
  \item \textbf{배열로의 일반화.}
    이 선호는 리액트의 의존성 배열에서만 나타나는 것이 아니다.
    \texttt{useEffect}나 \texttt{subscribe}처럼 함수와 배열을 인자로 전달하는 상황이 아니라 독립적인 배열을 완성할 때에도 똑같이 상태 변수를 더 선호한다.
    따라서 모델이 상태 변수를 선호하는 현상은 함수--배열 패턴에서만 일어나는 것이 아니라 모델이 배열을 채우는 일반적 경향에 가깝다.
\end{itemize}

\section{선행 연구}
\label{sec:related}

\paragraph{코드 모델의 표상 분석.}
코드 모델이 코드의 어떤 정보를 학습하고 표상하는지를 분석한 연구는 여러 갈래로 진행되어 왔다.
\citet{troshin2022probing}은 사전 학습된 코드 모델의 표상이 문법 구조, \alt{이름}{identifier} 및 \alt{이름공간}{namespace}, 데이터 흐름과 같은 정보를 얼마나 담고 있는지 분석하였다.
\citet{ma2024unveiling}은 이를 더 체계적으로 확장하여, 문법 구조, 제어 흐름 및 의존성, 데이터 의존성과 같은 코드 구조가 여러 코드 모델의 표상에 얼마나 반영되는지 살펴보았다.
\citet{maveli2025functional}은 \alt{기능적 동등성}{functional equivalence} 평가를 통해, 코드 모델이 표면적 구문을 넘어 프로그램 의미를 얼마나 포착하는지 물었다.
이들 연구는 코드 모델이 코드의 문법적·의미적 정보를 일부 표상한다는 배경을 제공하지만, 대체로 코드 전체나 코드 쌍, 전반적인 모델 표상처럼 거시적인 관점에서 분석을 진행하였다는 한계가 있다.

\paragraph{코드 모델의 표면 단서 민감성.}
한편 코드 모델의 예측이 프로그램의 의미를 안정적으로 따르지 않는다는 점을 보인 연구도 있다.
\citet{yefet2020adversarial}은 작은 코드 변형만으로도 모델의 예측이 흔들릴 수 있음을 보였고, \citet{wang2024naming}은 변수, 함수, 메서드 등의 이름이 코드 분석 과제의 성능에 상당한 영향을 줄 수 있음을 보였다.
\citet{orvalho2025robust}은 의미를 보존하는 변형 앞에서도 대형 언어 모델의 코드 이해와 예측이 불안정해질 수 있음을 밝혀냈다.
이들 연구는 코드 모델의 다음 토큰 예측이 코드의 표면적인 단서에 영향을 받을 가능성이 있음을 시사한다.

\paragraph{}
이처럼 기존 연구는 주로 완성된 코드의 정답 여부, 변형에 대한 강건성, 또는 모델 표상의 전반적 성질을 다루었다.
반면 이 연구는 모델의 행동적 선호를 국소적으로 관찰, 측정한다.
우리는 주어진 코드 조각의 다음에 이어질 토큰을 모델이 예측하도록 하고, 후보 토큰의 로짓 차이를 직접 비교하였다.
이를 통해 모델이 주어진 코드에서 어떤 행동을 취하는지, 또 어떤 표면 단서가 그 행동을 유발하는지를 밝히고자 한다.

\section{리액트 훅}
\label{sec:react-hooks}

\alt{리액트}{React}는 웹 기반 \alt{사용자 인터페이스}{UI}를 개발하기 위한 자바스크립트 라이브러리이다.
리액트로 만들어진 UI는 \alt{컴포넌트}{component}의 트리로 구성되며, 각 컴포넌트는 화면의 한 조각과 그 동작을 선언적으로 기술하는 명세로, \alt{부수효과}{side-effect}를 일으키지 않는 순수 함수의 형태이다.
리액트 런타임은 이 명세를 해석하기 위해 컴포넌트를 직접 호출하고, 컴포넌트는 화면에 무엇이 나타나야 하는지를 나타내는 자료구조를 반환한다.
이상적으로는, 같은 입력으로 호출된 컴포넌트는 같은 사용자 인터페이스를 반환해야 하지만, 실제 프로그램에서는 상태 변화나 사용자 상호작용, 네트워크 요청과 같은 부수효과도 함께 다루어야 한다\citep{lee2025react}.

\alt{훅}{Hook}은 함수형 컴포넌트 안에서 이러한 상태와 부수효과를 다룰 수 있게 해 주는 함수이다.
훅은 겉으로는 일반적인 함수 호출처럼 보이지만, 리액트 런타임과 연결되어 컴포넌트의 상태를 관리하거나 렌더링 이후 실행될 효과를 등록한다.
훅에는 여러 종류가 있지만, 이 연구에서는 \texttt{useState}와 \texttt{useEffect}, \texttt{useLayoutEffect}를 위주로 살펴본다.

\paragraph{\texttt{useState}.}
\texttt{useState}는 컴포넌트 안에서 상태를 관리하기 위한 훅이다.
\texttt{useState}는 초기 상태 $x_0$을 인자로 받아, 현재 상태를 나타내는 값 $x$와 그 상태를 갱신하는 함수 $x_\texttt{set}$의 순서쌍을 반환한다\citep{met2026stt}.
주로 아래와 같이 구조 분해 할당 문법과 함께 쓰인다.

% tex-fmt: off
\begin{reactcode}
const [s, setS] = useState(0);
\end{reactcode}
% tex-fmt: on

서론의 예시에서 본 \texttt{s}가 바로 이런 방식으로 선언된 상태 변수이고, 이에 반해 \texttt{const~c~=~0}처럼 일반적인 방식으로 선언된 \texttt{c}는 상수다.
이처럼 한 변수를 정의할 때는 \texttt{useState} 훅과 구조 분해 할당 문법, 갱신 함수의 존재를 동반하는 상태의 형태로도 선언할 수 있고, 일반 자바스크립트의 예약어인 \texttt{const}를 사용하여 상수의 형태로도 정의할 수 있다.
이 글에서는, 변수가 표면적으로 어떤 형태로 정의되었는지를 두고 `선언 형태'라고 부르기로 하자.

\paragraph{\texttt{useEffect}.}
\texttt{useEffect}는 특정 시점에 실행될 부수효과를 등록하기 위한 훅으로, 네트워크 요청이나 외부 시스템과의 동기화처럼 컴포넌트의 순수한 반환값만으로는 표현하기 어려운 동작을 기술하는 데 쓰인다.
\texttt{useEffect} 훅은 \alt{콜백 함수}{callback}~$f$와 배열 $[\Overline{e}]$를 인자로 받으며, 화면이 \alt{그려지거나}{rendered} 의존성 배열에 있는 값이 변화할 때마다 리액트 런타임이 함수 $f$를 실행한다\citep{met2026sync}.
이때 $f$를 효과 함수, $[\Overline{e}]$를 \alt{의존성 배열}{dependency array}이라고 부르자.

의존성 배열은 부수효과가 어떤 값의 변화에 반응해야 하는지를 표시하는 자리이므로, \citep{met2026sep}에서는 효과 함수 내부에서 참조되는 반응형 값만을 의존성 배열에 포함할 것을 권고한다.
반응형 값이란 컴포넌트 안에서 선언된 상태 등 프로그램을 실행하는 과정에서 변화할 수 있는 값을 의미한다.
반대 개념으로서 상수와 같은 비반응형 값이 있으며, 프로그램이 실행되는 동안 변화하지 않는 값을 이른다.

아래는 상태 \texttt{s}가 바뀌거나 화면이 그려질 때마다 \texttt{s}의 값을 콘솔에 출력하도록 하는 코드이다.

% tex-fmt: off
\begin{reactcode}
useEffect(() => {
  console.log(s);
}, [s])
\end{reactcode}
% tex-fmt: on

이처럼 \texttt{useEffect} 훅은 콜백 함수와 배열 인자를 외부 함수가 감싸는 형태로 쓰인다.
이 글에서는 이 문법 패턴을 `함수--배열 구조'라고 부른다.

\paragraph{\texttt{useLayoutEffect}.}
\texttt{useLayoutEffect}는 \texttt{useEffect}와 마찬가지로 부수효과를 등록하기 위한 훅이며, 효과 함수와 의존성 배열을 인자로 받는다는 점에서 동일한 함수--배열 구조를 가진다.

차이는 효과가 실행되는 시점에 있다.
\texttt{useEffect}에 등록된 효과는 일반적으로 화면이 그려진 뒤 비동기적으로 실행되는 반면, \texttt{useLayoutEffect}에 등록된 효과는 브라우저가 화면을 다시 그리기 전에 동기적으로 실행된다\citep{met2026lay}.
따라서 두 훅은 실행 시점과 사용 목적에서는 구별되지만, 이 연구의 관점에서는 모두 콜백 함수와 의존성 배열을 함께 받는 리액트 훅이라는 공통점을 가진다.

이 연구는 \texttt{useEffect}뿐 아니라 \texttt{useLayoutEffect} 조건도 함께 살펴봄으로써, 배열 첫 항목에 대한 모델의 선호가 특정 함수 이름 \texttt{useEffect}에만 묶인 것인지, 아니면 다른 \texttt{useEffect} 계열 훅에서도 유지되는지를 확인한다.

\bibliography{custom}
\end{document}
